\documentclass[11pt]{article}
\usepackage[a4paper,margin=27mm]{geometry}
\usepackage{amsmath,amssymb,amsthm,longtable,array}
\usepackage[hidelinks]{hyperref}
\newtheorem{theorem}{Theorem}[section]
\newtheorem{lemma}[theorem]{Lemma}
\newtheorem{proposition}[theorem]{Proposition}
\newtheorem{corollary}[theorem]{Corollary}
\newcommand{\gen}[1]{\langle #1\rangle}
\newcommand{\Aut}{\operatorname{Aut}}
\title{Coset residues and chamber orbits}
\author{Human Compression project}
\date{}
\begin{document}
\maketitle
\begin{abstract}
We give a complete account of the valid residue arguments and the
necessary corrections in Leemans--Tranchida's paper on residual
connectedness. The rank-three result and a criterion using the
specified group's action are proved. The published general criterion,
which uses chirality under the full automorphism group, has a previously
verified counterexample. Two smaller examples dispose of auxiliary
assertions whose wording omits firmness or a rank restriction. The
counterexamples, their scope and the corrected theorems are kept
distinct. This manuscript assembles prior project proofs and the
accepted corrections; it makes no additional conjecture or measurement
claim.
\end{abstract}

\section{Source and conventions}
The source is the final journal edition of Leemans--Tranchida
\cite{LT}, pages 491--499. Its principal Theorem 1.1 asserts residual
connectedness for the standard coset geometry of a $C^+$-group if that
geometry is chiral. Chirality is defined using the \emph{full}
type-preserving automorphism group. That group need not be the
specified $C^+$-group. The accepted five-page manuscript \cite{CE}
gives an explicit counterexample with specified group
$\operatorname{PSL}_2(13)$ and full group
$\operatorname{PSL}_2(13)\times C_2$. It is reproduced in the complete
volume with its proof unchanged.

We retain the source's definition of a geometry: every flag extends
to a chamber. Firmness means that every panel has at least two
completions; thinness means exactly two. Residual connectedness asks
that every residue of rank at least two, including the whole geometry,
have a connected incidence graph. None of these extra conditions is
implicit in the word geometry. In particular, a rank-one residue with
two objects has a disconnected incidence graph.

Let $I$ be a finite type set. We use left cosets $gH_i$ and incidence
by nonempty intersection, retaining the type as part of an object.
The source uses right cosets. Inversion sends $H_i g$ to $g^{-1}H_i$
and is an isomorphism of the two incidence systems. Thus this
convention changes no flags, residues or transitivity assertions.
Write $H_J=\bigcap_{j\in J}H_j$, with $H_\varnothing=G$.
All general proofs below permit infinite groups.

The coset construction, the $C^+$ intersection condition and the
canonical neighbouring chambers are prior results of the source and
its cited predecessors. Sections 2--5 rework the existing project
arguments \cite{IC,CC}. The chamber-component explanation in Section 6
is also retained from \cite{CC}. The principal counterexample
\cite{CE}, the rank boundary \cite{PB}, and the seven-object example
\cite{CF} have distinct reviewed proofs. Classical orbit arguments are
not claimed as new. No new Lean proof is asserted by this assembly.

\section{Flags and the internal residue map}
The base objects $H_i$ form a chamber. Each object is a translate of
one of them. If $gH_i$ and $hH_j$ are incident, choose an element $x$
of their intersection. Their pair is then $x\{H_i,H_j\}$. This proves
the object and pair transitivity in the cited Tits construction, and
also proves that every flag of rank at most two extends to a chamber.

Fix a proper subset $J\subset I$, and let $F_J=\{H_j:j\in J\}$.
For $i\notin J$, put
\[
 S_{J,i}=\bigcap_{j\in J}H_jH_i.
\]
For $J=\varnothing$ the intersection is $G$. The type-$i$ objects
of the residue of $F_J$ are exactly the cosets $gH_i$ with
$g\in S_{J,i}$.

\begin{lemma}\label{internal}
The maps
\[
 h(H_J\cap H_i)\longmapsto hH_i\qquad(h\in H_J, i\notin J)
 \tag{1}
\]
define a homomorphism from the internal coset incidence system
$\Gamma(H_J;(H_J\cap H_i)_{i\notin J})$ into the residue of $F_J$.
It is injective on objects. It is surjective on objects if and only if
$S_{J,i}=H_JH_i$ for every $i\notin J$.
\end{lemma}
\begin{proof}
For $h,k\in H_J$, equality of $hH_i$ and $kH_i$ is equivalent to
$k^{-1}h\in H_i$, and hence to $k^{-1}h\in H_J\cap H_i$.
This proves both well-definedness and injectivity. A common element
of two internal cosets is a common element of the corresponding
external cosets, proving preservation of incidence. The image at
type $i$ consists of the cosets represented by $H_JH_i$. The full
residue at that type is represented by $S_{J,i}$, and both sets are
unions of left cosets of $H_i$. Equality is therefore exactly the
surjectivity condition.
\end{proof}

A bijective incidence homomorphism need not be an isomorphism: its
inverse may fail to preserve incidence. This distinction in the source's
Remark 4.4 is essential. For a singleton $J=\{j\}$ the displayed
surjectivity condition is automatic, but no assertion about the inverse
follows from this alone.

\begin{lemma}\label{extension}
Suppose the coset system is a geometry. The action of $G$ is
transitive on flags of type $J\cup\{i\}$ if
and only if it is transitive on flags of type $J$ and
$S_{J,i}=H_JH_i$.
\end{lemma}
\begin{proof}
Transitivity on the larger type implies transitivity on the smaller
type, because every smaller flag has an extension in a geometry.
For $g\in S_{J,i}$, the flag $F_J\cup\{gH_i\}$ is of the larger
type. An element carrying the corresponding base flag to it fixes
every member of $F_J$, so lies in $H_J$. Hence $gH_i=hH_i$ for
some $h\in H_J$ and $g\in H_JH_i$. The reverse inclusion always
holds.

Conversely, normalize the $J$-part of any flag of the larger type
to $F_J$. Its last object is $gH_i$ with $g\in S_{J,i}$. Write
$g=hu$ with $h\in H_J$ and $u\in H_i$. Multiplication by $h^{-1}$
fixes $F_J$ and takes the last object to $H_i$. Thus every such flag
is a translate of the base flag.
\end{proof}

This is the source's Lemma 4.5. If the maps (1) are surjective for
every $J$, Lemma~\ref{extension}, by induction on the size of a flag,
gives flag transitivity. In that case (1) is an isomorphism: any
incident pair in a residue, together with $F_J$, is a flag and can
be normalized by an element of $H_J$. That element is a common
representative in the internal model. This proves the final implication
in the cited source Lemma 4.3 for the maps used here.

Surjectivity for just one $J$ does \emph{not} supply transitivity on
$J$. This is the missing condition in Proposition 4.6. The retained
counterexample is $G=C_2^2$, with its three order-two subgroups
$H_0,H_1,H_2$ and $H_3=G$. Its geometry is the cone over the complete
tripartite graph $K_{2,2,2}$, with seven objects and eight chambers.
For $J=\{0,1,2\}$ the base residue and its internal coset model each
have one object, so (1) is surjective. But a group of order four
cannot be transitive on the eight $J$-flags. This is a counterexample
to the stated equivalence, not to a $C^+$ assertion.

\section{Connectivity and the specified action}
\begin{lemma}\label{connected}
Let $K_i$ be subgroups of a group $K$ and let the index set have at
least two elements. The incidence graph of their coset system is
connected if and only if $K=\gen{K_i:i\in I}$.
\end{lemma}
\begin{proof}
The base cosets form a clique. Each $K_i$ fixes one of its vertices,
and therefore preserves the connected component of that clique.
If the subgroups generate $K$, every translate of every base vertex
lies in this component, proving connectedness. Conversely, a path
starting at a base vertex can be followed by successively choosing
elements in the intersections of its consecutive cosets. Each step
changes the representative by multiplication by an element of one
of the $K_i$. Every vertex reached therefore has a representative in
$\gen{K_i:i\in I}$. If all vertices are reached, each coset of any
one $K_i$ has such a representative, so their union gives
$K=\gen{K_i:i\in I}$.
\end{proof}

\begin{proposition}\label{properconnected}
Suppose $G$ is transitive on flags of every proper type in a coset
geometry. If $|I\setminus J|\geq2$ and
\[
 H_J=\gen{H_J\cap H_i:i\in I\setminus J},                 \tag{2}
\]
then every residue of a flag of type $J$ is connected.
\end{proposition}
\begin{proof}
Normalize the flag to $F_J$. For every $i\notin J$, the set
$J\cup\{i\}$ is still proper, so Lemma~\ref{extension} and
Lemma~\ref{internal} make the internal model bijective on objects.
Equation (2) and Lemma~\ref{connected} make that internal graph
connected. Its image is a connected spanning subgraph of the actual
residue. Any extra incidence edges preserve connectedness.
\end{proof}

This is the sufficient residue statement used in the source, with
its action and rank hypotheses explicit. In rank one, transitivity
on the enlarged type would mean transitivity on chambers, which
has not been assumed. The rank-one boundary of its printed
Proposition 4.7 is exhibited in \cite{CF}. The source's preceding
paragraph already excludes these residues from the residual-connectedness
question.

\section{Canonical chambers of a marked group}
From now on let $|I|=r\geq3$, with a distinguished element $0$.
Let $G$ be generated by distinct nonidentity elements $\alpha_i$
for $i\ne0$, and put $\alpha_0=1$. Set
\[
 a_{ij}=\alpha_i^{-1}\alpha_j,\qquad
 P_S=\gen{a_{ij}:i,j\in S},\qquad H_i=P_{I\setminus\{i\}}.
\]
Assume the $C^+$ intersection condition
\[
 P_S\cap P_T=P_{S\cap T}\qquad(|S|,|T|\geq2).             \tag{3}
\]
For an empty or singleton set $S$, the group $P_S$ is trivial, so
(3) holds also when either index set has size less than two.
Repeated intersection therefore gives
\[
                  H_J=P_{I\setminus J}.                  \tag{4}
\]
For $i>0$, $H_i$ is generated by the $\alpha_j$ with $j\ne i$;
for $i=0$, it is generated by $\alpha_1^{-1}\alpha_j$ for $j\geq2$.
Thus these are precisely the maximal parabolics in Construction 3.2.

\begin{lemma}\label{neighbours}
Let $C=\{H_i:i\in I\}$. For $i\in I$ and any $j\ne i$, put
\[
 D_i=(C\setminus\{H_i\})\cup\{a_{ij}H_i\}.
\]
This does not depend on $j$. It is a chamber distinct from $C$,
adjacent to it at type $i$. Moreover
\[
 a_{ji}D_i=D_j,\qquad D_i=\alpha_i^{-1}D_0,
 \qquad D_0\text{ is $i$-adjacent to }\alpha_i C.          \tag{5}
\]
The two $G$-orbits $GC$ and $GD_0$ are disjoint.
\end{lemma}
\begin{proof}
For $j,k\ne i$, the quotient $a_{ij}^{-1}a_{ik}=a_{jk}$ belongs
to $H_i$, proving independence of the choice. The coset differs
from $H_i$ because $a_{ij}\ne1$ and
$P_{\{i,j\}}\cap H_i=P_{\{j\}}=1$.
If $\ell\ne i,j$, its representative $a_{ij}$ belongs to $H_\ell$.
To obtain incidence with $H_j$, choose $k\ne i,j$; this is where
$r\geq3$ is needed. The element $a_{ik}=a_{ij}a_{jk}$ lies in both
$a_{ij}H_i$ and $H_j$. All unchanged objects meet at the identity,
so $D_i$ is a chamber.

For the first identity in (5), at type $i$ the two multipliers
cancel; at type $j$ the multiplier produces the altered coset;
at every other type it lies in that parabolic. Taking one index
to be zero gives the second identity. Translating the adjacency
of $C,D_i$ by $\alpha_i$ gives the last assertion in (5).
If $gC=D_0$, then $g$ fixes the base panel missing type zero.
Its stabilizer is $H_{I\setminus\{0\}}=P_{\{0\}}=1$ by (4).
This would give $D_0=C$, a contradiction.
\end{proof}

The lemma is the valid rank-at-least-three form of source Lemma 4.8.
It describes neighbours of these canonical chambers. It does not
normalize an arbitrary chamber to $C$.

\begin{proposition}\label{firmchiral}
If the marked coset system is a geometry and its full type-preserving
group $A$ has two alternating chamber orbits, then it is thin.
For an arbitrary firm geometry with such a group action, $A$ is
transitive on flags of every proper type, and residual connectedness
can be tested on the residues of a single chamber.
\end{proposition}
\begin{proof}
In the marked case, $C$ and $D_0$ are adjacent and therefore lie
in the two different $A$-orbits. Lemma~\ref{neighbours} supplies
a neighbour of every type for both orbit representatives. Transport
by $A$ supplies one for every chamber. Every panel has at least
two completions. It cannot have three: three chambers completing
one panel are pairwise adjacent, and three pairwise adjacent
chambers cannot lie alternately in just two orbits. Thus the
geometry is thin.

For the general firm case, let $F$ be a flag of proper type and
choose a chamber containing it. At a type missing from $F$,
firmness gives an adjacent chamber still containing $F$. These
two chambers lie in the two different orbits. Given another flag
of the same type, extend it to a chamber and compare with the
first or second chamber as required. The resulting automorphism
maps the whole typed flag to the other one. Thus $A$ is transitive
on each proper flag type, and their residues are isomorphic.
Every proper type occurs in a chosen chamber; this proves the
residue test, including all residues of rank at least two.
\end{proof}

This proves the useful firm versions of source Theorem 4.1 and
Proposition 4.2. It does not prove their unrestricted wording:
the seven-object example in \cite{CF} satisfies the source's
literal chirality definition and refutes both. Nor does this
proposition give proper-flag transitivity for the subgroup $G$.

\section{The valid residual-connectedness criterion}
\begin{theorem}\label{criterion}
For the nondegenerate $C^+$ marking above, suppose its coset system
is a geometry and $G$ is transitive on flags of every proper type.
Then the geometry is firm and residually connected.
\end{theorem}
\begin{proof}
Each panel can be normalized to a base panel. Its two canonical
completions from Lemma~\ref{neighbours} give firmness.
Let $F_J$ be a base flag, put $S=I\setminus J$, and assume
$|S|\geq2$. As in Proposition~\ref{properconnected}, the internal
residue model is bijective on objects and preserves incidence.

If $|S|\geq3$, equation (4) identifies its ambient group with
$P_S$ and its parabolics with $P_{S\setminus\{i\}}$ for $i\in S$.
Every generator $a_{uv}$ of $P_S$ belongs to one of those parabolics:
choose $i\in S\setminus\{u,v\}$. They therefore generate $P_S$,
and Proposition~\ref{properconnected} proves connectedness.

If $S=\{u,v\}$, write $a=a_{uv}$, so $H_J=\gen{a}$.
Its two internal parabolics are trivial. The residue has two
copies of this cyclic group as its objects. Translates of $C$
give edges from $a^kH_u$ to $a^kH_v$. Translates of $D_v$
give edges from $a^kH_u$ to $a^{k-1}H_v$.
Here $a\ne1$, and distinct powers represent distinct cosets
until the order of $a$ is reached, by (4). These two matchings
form a cycle of length $2|a|$ if $a$ has finite order, and a
doubly infinite path otherwise. Their union is connected and
spans the residue. Possible additional edges cannot destroy
connectedness. This also explains the order-three diagram of
source Example 4.10; reversing both cyclic coordinates changes
the sign of the shift without changing the argument.

Finally every flag of proper type can be normalized to $F_J$,
so the conclusion holds for all residues, including the empty
flag's residue.
\end{proof}

\begin{corollary}\label{rankthree}
Every rank-three system from a nondegenerate $C^+$ marking is a
firm, residually connected geometry.
\end{corollary}
\begin{proof}
Every incident pair is a translate of a base pair and extends
to a chamber. Objects and the empty flag extend too; any flag
of size three is already a chamber. Thus it is a geometry.
Every proper flag has rank at most two, so the coset action is
transitive on its type. Theorem~\ref{criterion} applies.
\end{proof}
This is source Proposition 4.11. It requires neither chirality
nor a full automorphism computation.

For rank two, the nontrivial cyclic marking has both parabolics
trivial and gives disjoint edges. The full written proof and
literal check in \cite{PB} refute the unrestricted Lemma 4.8
and Corollary 4.9. Imposing rank at least four in a proposed
replacement for Corollary 4.9 produces a different assertion,
unproved here. A false universal statement is not recovered
by silently changing its hypotheses, and its counterexample
does not automatically create a new source obligation for
every such replacement.

\section{Why the two acting groups differ}
The following retained argument pinpoints the principal error.
It adds a structural explanation to the explicit counterexample.

\begin{theorem}\label{component}
If the marked coset system is a thin geometry, then
$\Omega=GC\mathbin{\dot\cup}GD_0$ is one gallery component.
It contains every object and every incident pair. The stabilizer
of $C$ in the full type-preserving group $A$ is trivial.
If $A$ is chiral, then $\operatorname{Stab}_A(\Omega)=G$ and
the gallery components are indexed by $A/G$.
\end{theorem}
\begin{proof}
Thinness makes the $i$-neighbour of $gC$ equal to $gD_i$,
which lies in $GD_0$. The $i$-neighbour of $gD_0$ is
$g\alpha_iC$. Thus $\Omega$ is closed under adjacency.
The two-step galleries $C,D_0,\alpha_iC$, their reversals,
and their translates join $C$ to every $gC$, because the
$\alpha_i$ generate $G$. Each $gD_0$ is adjacent to $gC$.
Consequently $\Omega$ is a connected component.

Every object occurs in some $gC$. For an incident pair,
an element of the intersection of the two cosets is a common
representative, placing the pair in a translate of $C$.
An automorphism fixing $C$ fixes its neighbours, since each
panel has a unique other completion. Repetition along galleries
fixes $\Omega$ pointwise. Since $\Omega$ contains every object,
the automorphism is the identity. The coset action itself is
faithful because its kernel fixes $C$ and lies in
$\bigcap_iH_i=P_\varnothing=1$.

Suppose now $A$ is chiral. If $f\in A$ preserves $\Omega$,
then $fC$ lies in $GC$ or $GD_0$. The second possibility would
identify the two $A$-orbits of the adjacent chambers $C,D_0$.
Hence $fC=gC$ for some $g\in G$. The element $g^{-1}f$
fixes $C$ and is the identity, so $f=g$. Conversely $G$
preserves $\Omega$. Every chamber is an $A$-translate of
$C$ or $D_0$, so every component is an $A$-translate of
$\Omega$. The orbit-stabilizer bijection gives $A/G$.
\end{proof}

The theorem does not assert that $A/G$ has one element.
The accepted example \cite{CE} has $|G|=1092$, $|A|=2184$,
two gallery components, and a disconnected rank-two residue.
Its full-group chirality supplies the hypotheses of the published
Theorem 1.1 and disproves its conclusion. The switch from
proper-flag transitivity for $A$ to that property for $G$ is
therefore a substantive error, separate from the smaller missing
rank and firmness conditions.

\section{Disposition of the source}
Figure~\ref{seven} illustrates the smaller counterexample from
\cite{CF}. Its chambers are the four triangles
$C_j=\{c_j,x_j,y_j\}$ and $D_j=\{d,x_j,y_j\}$, for $j=1,2$.
The full typed group exchanges the two halves. Its chamber orbits
are $\{C_1,C_2\}$ and $\{D_1,D_2\}$, and the only adjacencies
join $C_j$ to $D_j$. Yet the panel $\{c_j,x_j\}$ has only one
completion. This is the missing-firmness issue; it is separate
from the principal $\operatorname{PSL}_2(13)$ counterexample.

\begin{figure}[ht]
\centering
\setlength{\unitlength}{.65mm}
\begin{picture}(220,78)
\put(10,35){\line(2,1){50}}
\put(10,35){\line(2,-1){50}}
\put(60,10){\line(0,1){50}}
\put(60,60){\line(2,-1){50}}
\put(60,10){\line(2,1){50}}
\put(110,35){\line(2,1){50}}
\put(110,35){\line(2,-1){50}}
\put(160,10){\line(0,1){50}}
\put(160,60){\line(2,-1){50}}
\put(160,10){\line(2,1){50}}
\put(10,35){\circle*{2}}
\put(60,60){\circle*{2}}
\put(60,10){\circle*{2}}
\put(110,35){\circle*{2}}
\put(160,60){\circle*{2}}
\put(160,10){\circle*{2}}
\put(210,35){\circle*{2}}
\put(0,33){$c_1$}
\put(57,65){$x_1$}
\put(57,0){$y_1$}
\put(108,40){$d$}
\put(157,65){$x_2$}
\put(157,0){$y_2$}
\put(213,33){$c_2$}
\end{picture}
\caption{All incidences of the seven-object example. The type classes
are $\{d,c_1,c_2\}$, $\{x_1,x_2\}$ and $\{y_1,y_2\}$.}
\label{seven}
\end{figure}

The complete source ledger accompanies this manuscript. The
following table records the mathematical outcomes; it is not
a claim that false statements have been reproved.
\begin{longtable}{p{.23\linewidth}p{.70\linewidth}}
\textbf{Source assertion}&\textbf{Disposition and proof}\\\hline
Theorem 1.1 & Refuted in its literal full-group scope by the
accepted five-page counterexample \cite{CE}. The valid specified-action
criterion is Theorem~\ref{criterion}.\\[3pt]
Proposition 2.1 & Cited Tits construction; object and pair
normalization proved in Section 2.\\[3pt]
Theorem 2.2 & Cited flag-transitive residue theorem from the
source's reference [1]. Its sufficient part needed here is proved
in Section 3; the whole cited monograph is not an additional intake.\\[3pt]
Theorems 3.1, 3.3;
Corollary 3.4 & Explicitly borrowed C-group and $C^+$ correspondence
results from the source's reference [6], \emph{Highly symmetric
hypertopes}. Their attribution is retained; the present corrections
do not refute them.\\[3pt]
Construction 3.2 & The parabolics are identified with
$P_{I\setminus\{i\}}$ in Section 4; left and right cosets are related
by the explicit inversion map in Section 1.\\[3pt]
Chirality implies thinness;
Theorem 4.1;
Proposition 4.2 & False without firmness under the stated geometry
definition; the seven-object example \cite{CF} supplies the full
group and all residues explicitly. The firm normalization statements
are proved in Proposition~\ref{firmchiral}.\\[3pt]
Lemma 4.3;
Remark 4.4 & The cited internal map, its object-surjectivity criterion,
and its distinction from an isomorphism are proved in Section 2 for
the interfaces used here.\\[3pt]
Lemma 4.5 & Recovered as Lemma~\ref{extension}, with the smaller
flag-type transitivity present.\\[3pt]
Proposition 4.6 & Refuted by the retained $C_2^2$ example in Section 2.
Adding smaller-type transitivity repairs the equivalence.\\[3pt]
Proposition 4.7 & The displayed rank-one boundary is false, as shown
in \cite{CF}; the preceding source paragraph focuses on rank at
least two. Proposition~\ref{properconnected} states the valid
rank and action hypotheses.\\[3pt]
Lemma 4.8;
Corollary 4.9 & Refuted in unrestricted rank by \cite{PB}. The
canonical-neighbour lemma for rank at least three is
Lemma~\ref{neighbours}. General firmness in rank at least four
without further action hypotheses remains unproved.\\[3pt]
Example 4.10;
Figure 1 & The cyclic two-matching construction is proved for all
cyclic orders in Theorem~\ref{criterion}; the source's order-three
illustration is a special case.\\[3pt]
Proposition 4.11 & Recovered as Corollary~\ref{rankthree}.\\
\end{longtable}

No named open question is solved or added by this disposition.
The principal refutation, valid source results, missing hypotheses,
and useful unproved variants are separate outcomes. Existing
computational certificates are reused without repeating their
acquisition. The exact finite checks in \cite{PB,CF} verify their
stated examples, not the arbitrary-group proofs above. Compression
measurements and further Lean formalization are deferred.

\begin{thebibliography}{9}
\bibitem{LT} D.~Leemans and P.~Tranchida,
\emph{On residual connectedness in chiral geometries},
Algebraic Combinatorics \textbf{4} (2021), 491--499,
\href{https://doi.org/10.5802/alco.162}{doi:10.5802/alco.162}.
\bibitem{IC} Human Compression project,
\emph{Incidence results derived from relation matrices},
retained proof bank, Section 1.
\bibitem{CC} Human Compression project,
\emph{Chamber components in marked coset geometries},
retained phase276 manuscript, Sections 1--5.
\bibitem{CE} Human Compression project,
\emph{A chiral coset geometry with disconnected residues},
accepted five-page manuscript and exact certificate.
\bibitem{PB} Human Compression project,
\emph{Panels in the coset construction for a $C^+$-group},
accepted three-page manuscript, with the rank boundary and the
retained finite panel checks.
\bibitem{CF} Human Compression project,
\emph{Chirality without firmness}, accepted two-page manuscript
and seven-object certificate.
\end{thebibliography}
\end{document}
